\section{Fresh corridors and a subcritical infinite cluster}
\label{sec:corridor-exploration}

Proposition~\ref{prop:moves} supplies paths with approximate coordinate
displacements $(1,1)$ and $(1,-1)$. We will use them to transfer a
connection through long corridors. The corridor axes have slopes $1/4$ and $-1/4$:
at each step, choosing the sign of the vertical move corrects the
path's displacement from its axis. The smaller slope leaves room for
this correction even when the entrance point is not prescribed.

The corridors form the bonds of an oriented square lattice, but their
vertex sets are coordinate preimages in $G$. All paths and revealed
edges remain in $G$. We will reveal an incoming corridor only after
estimating connections through the still fresh outgoing corridors.
This order will give a uniform conditional failure bound for each
tested lattice site.

\subsection{The deterministic layout}

The constant $C$ is the path-displacement bound from
Proposition~\ref{prop:moves}. Let $B$ be the initial vertical tolerance
for an entrance, and let $H$ be the tube half-width. Fix
\[
 C=100,\qquad \mu=\tfrac14,\qquad B=2C+10,\qquad H=3B.
\]
An endpoint rectangle of time half-width $J$ will contain every overlap
of incident tubes. Choose that width first, and then the corridor
length, by taking integers $J,L$ with
\[
 J>10(H/\mu+C),\qquad L>10(J+B),\qquad T=L+\tfrac12.
\]
For a site $v=(a,b)\in\mathbb N_0^2$, let
$c_v=((a+b)T,(a-b)\mu T)$. Its children are $(a+1,b)$ and
$(a,b+1)$; their displacement from $v$ is $(T,\sigma\mu T)$ with
$\sigma=1$ and $-1$, respectively. All regions below mean vertex
preimages under the coordinates $(t,z)$, with coordinates taken
relative to the indicated center. Define
\[
 Q_v=\{|t-c_v^{(1)}|\le J,\
          |z-c_v^{(2)}|\le\mu J+H\}.
\]
For a parent--child pair $v,x$ of sign $\sigma$, let $D_{vx}$ be the
induced subgraph on the union of $Q_v,Q_x$ and the tube
\[
 0\le t-c_v^{(1)}\le T,\qquad
 |z-c_v^{(2)}-\sigma\mu(t-c_v^{(1)})|\le H.
\]
Write $\mathcal E_{vx}$ for its edge set, and $\mathcal Q_v$ for
the edge set induced by $Q_v$.

\begin{lemma}[Corridor separation]\label{lem:corridors}
Distinct $Q_v$ are disjoint. Two distinct corridors with no common
endpoint have disjoint vertex sets; corridors sharing an endpoint $v$
intersect only in $Q_v$. A corridor misses every nonendpoint $Q_w$.
Consequently, common edges of incident corridors belong to
$\mathcal Q_v$ at their common endpoint.
\end{lemma}

\begin{proof}
Site layers are spaced by $T$ in the time coordinate; within a layer
their heights are spaced by $2\mu T$. The choices give
$T>2J$ and $\mu T>\mu J+H$, so the $Q$ regions are disjoint.
A tube near either endpoint layer has height differing from that
endpoint by at most $\mu J+H$ when its time differs by at most $J$.
It therefore misses every other $Q$ in that layer; other layers are
excluded by their time coordinates.

Tubes in nonadjacent layer intervals have disjoint time interiors.
At a common layer plane, distinct endpoint heights differ by
$2\mu T>2H$, so adjacent-layer tubes can meet only at a shared
endpoint, inside its $Q$. In the same layer interval, parallel tube
axes are separated by $2\mu T$. Oppositely sloped axes can approach
within $2H$ only when they share a start, within time $H/\mu$ of
that start, or share an end, within time $H/\mu$ of that end.
These intersections lie in the corresponding $Q$, since $J>H/\mu$.
Combining these tube facts with the $Q$ facts proves the vertex claims.
An edge in both induced subgraphs has both endpoints in their vertex
intersection, proving the last assertion. Edges joining different
corridors are never added merely because both their endpoints have
been revealed.
\end{proof}

\subsection{Relay regions and the order of scales}
\label{subsec:corridor-scales}

Inside a corridor we will place successive input regions near integer
times. Their vertical half-width decreases from $B$ by $1/2$ per step
until it reaches $2$. The local moves can make this correction because
their vertical displacement has magnitude one while the axis advances
only $\mu=1/4$. The final move will be stopped at the child mark,
halfway between the last two integer times.

For a vertex set $A$, its distance-$R$ collar is the ambient
neighborhood $\{x\in V(G):d_G(x,A)\le R\}$, where
$d_G(x,A)=\inf_{a\in A}d_G(x,a)$ and $\inf\varnothing=\infty$.
We need each input region's collar to contain many disjoint finite
seed trials. The collar may be wide in graph distance while narrow in
coordinates, because the edge-coordinate bound is chosen last.

Let $\varepsilon$ be the permitted conditional probability of a failed
lattice site. Choose it so small that
\begin{equation}\label{eq:directed-contour-series}
 \sum_{l\ge1}4^l\varepsilon^{l/4}<1.
\end{equation}
The threshold $\delta$ will test stored connection probabilities;
$\kappa$ will bound the loss at a single relay. Set
\[
 \delta=\varepsilon/4,\qquad
 \kappa=\frac{\varepsilon\delta}{8(L+1)},\qquad
 \lambda=\kappa/4,\qquad K=\log(4/\kappa).
\]
The transfer proof will use the auxiliary error $\lambda$ and trial
weight $K$. To keep accumulated endpoint errors below $1/100$, choose
$e>0$ with $(L+3)e<1/100$. Choose the failure probability of a seed
move to satisfy
\[
 0<\eta_0\le\frac{\kappa\lambda}{4(K+1)}.
\]
Proposition~\ref{prop:moves}, with $e,\eta_0$, first gives a seed
radius $m$. We next choose the number $k$ of possible contacts needed
for disjoint seed trials and the graph-distance collar width $R$.
With $D$ the maximum degree, take integers
\begin{equation}\label{eq:corridor-trial-scales}
 \begin{split}
 k&>V_+(4m+2)
       \left\lceil K/(\pc/2)^{m+1+V_+(m)}\right\rceil,\\
 R&>m+4(1-\pc)^{-Dk}/\kappa.
 \end{split}
\end{equation}
Now choose $\rho>0$ so that
\begin{equation}\label{eq:corridor-coordinate-step}
 (R+m+1)\rho<e/10,
\end{equation}
and obtain the coordinates and a common $p'\in[\pc/2,\pc)$ from
Proposition~\ref{prop:moves}, with edge increments at most $\rho$.
Thus $m$ is fixed before the trial count and collar width; both are
fixed before the coordinate scale and the parameter $p'$.

At site $v$ put
\[
 M_v=\{0\le t-c_v^{(1)}\le\rho,\
                   |z-c_v^{(2)}|\le B\}.
\]
In a corridor $v,x$ of sign $\sigma$, use the successive input regions
$B_0=M_v$ and, for $1\le i\le L$,
\begin{equation}\label{eq:relay-regions}
 B_i=\{|t-c_v^{(1)}-i|\le(i+1)e,\
 |z-c_v^{(2)}-\sigma\mu i|\le b_i\},
 \qquad b_i=\max(2,B-i/2).
\end{equation}
The last target is $B_{L+1}=M_x$.

Each distance-$R$ collar of $B_i$, $0\le i\le L$, is contained in
$D_{vx}$. From every vertex $y$ in that collar, one of the two seed
moves reaches the next region with probability at least $1-\eta_0$,
using paths in $D_{vx}$. We give the estimates, including the last
crossing. Suppose first that $0\le i<L$ and that $y$ is in the
collar of $B_i$. A distance-$R$ displacement changes either coordinate
by at most $R\rho$. Write
$r=z(y)-c_v^{(2)}-\sigma\mu i$ for the vertical error at the current
integer time. Choose a move sign $\tau\in\{-1,1\}$ with $\tau r\le0$,
either sign when $r=0$. Its target error relative to the next axis
point is $r+\tau-\sigma\mu$, up to $e/4$. With $b_0=B$, its magnitude
is at most
\[
 \max(1+\mu,b_i+R\rho-1+\mu)+e/4\le b_{i+1}.
\]
The inequality follows from $\mu=1/4$, $b_i\ge2$, and
$R\rho+e/4<1/4$. Time error increases by at most $R\rho+e/4<e$;
at $i=0$ the additional initial interval length $\rho$ also fits
the allowed $2e$.

Since $b_L=2$, a move from the collar of $B_L$ starts below time $T$
and ends above it. The starting vertex may be anywhere in $B(y,m)$,
so its additional time error is $m\rho$, paid for by
\eqref{eq:corridor-coordinate-step}. Along the path, the first vertex
at or above $T$ has time in $[T,T+\rho]$. Its height differs from
the child center by at most $2+R\rho+C+\mu/2<B$, so it lies in $M_x$.

For every move, while its time is between the endpoint times, its
deviation from the tube axis is at most
\[
 B+R\rho+C+\mu\bigl((L+1)e+\rho+R\rho+C\bigr)<H.
\]
Outside that time interval its time is within $C+1$ of an endpoint
and its height lies in that endpoint's $Q$. This proves containment
of the paths. Collar containment follows from the same estimates
without the $C$ terms. The induced corridor therefore contains every
original edge of each indicated path, as required for the seed-move
probability to hold inside the corridor.

\begin{figure}[htbp]
\centering
\begin{tikzpicture}[x=.9cm,y=.9cm,>=Stealth]
 \foreach \ax/\ay/\bx/\by in {0/0/5/2.5,0/0/5/-2.5,5/2.5/10/0,5/-2.5/10/0}{
  \filldraw[fill=blue!8,draw=blue!45!black,line width=.5pt]
    (\ax,{\ay+.2})--(\bx,{\by+.2})--(\bx,{\by-.2})--(\ax,{\ay-.2})--cycle;
  \draw[->,blue!55!black,line width=.7pt,shorten <=9pt,shorten >=9pt]
    (\ax,\ay)--(\bx,\by);
 }
 \foreach \ax/\ay in {0/0,5/2.5,5/-2.5,10/0}{
  \filldraw[fill=orange!10,draw=orange!65!black]
    ({\ax-.42},{\ay-.52}) rectangle ({\ax+.42},{\ay+.52});
  \filldraw[fill=green!30!black,draw=green!30!black]
    (\ax,{\ay-.3}) rectangle ({\ax+.04},{\ay+.3});
 }
 \node[left=7pt] at (-.42,0) {$Q_{(0,0)}$};
 \node[above=6pt] at (5,3.02) {$Q_{(1,0)}$};
 \node[below=6pt] at (5,-3.02) {$Q_{(0,1)}$};
 \node[right=7pt] at (10.42,0) {$Q_{(1,1)}$};
 \node[above=4pt,rotate=26.565] at (2.5,1.48) {$D_{(0,0),(1,0)}$};
 \draw[->,black!65] (7,-2.15)--(9.97,-.28);
 \node[below,align=center] at (7,-2.18) {mark $M_{(1,1)}$};
\end{tikzpicture}
\caption{A diamond in the corridor layout. Corridors intersect only
inside the rectangle at a shared endpoint. The narrow green regions
are the marks $M_v$. Revealing an incoming corridor reveals the
endpoint rectangle, while the parts of the outgoing corridors outside
that rectangle remain fresh. Widths and the two coordinate scales are exaggerated;
the separation inequalities are proved in Lemma~\ref{lem:corridors}.}
\label{fig:corridors}
\end{figure}


\subsection{A transfer estimate in a fresh collar}

The exterior-contact, independent-seed, and shell-conditioned target
construction adapts Kozma and Nitzan's Lemma~10
\cite[Section~4]{KozmaNitzan2024}. We prove the version needed for fresh
metric collars in the original graph, using Theorem~\ref{thm:joint-gluing}.

The next argument applies to a countable product graph consisting of
original edges of $G$ and one synthetic source $*$ wired to finitely
many vertices. Edges outside the designated fresh region may have
arbitrary independent probabilities, including zero and one.

\begin{proposition}[Fresh-collar transfer]\label{prop:transfer}
Let $G$ and the parameters $m,k,R,K,\eta_0,\lambda,\kappa,p'$ satisfy
the choices in Subsection~\ref{subsec:corridor-scales}. Consider a
countable product graph whose edges are original edges of $G$,
together with a synthetic source $*$ joined by deterministic open
edges to finitely many vertices. Original edges have arbitrary
independent probabilities except where specified below.
Let $T_*$ be a fixed vertex set in this product graph.
Let $B_*$ be a vertex set whose distance-$R$ collar is contained in a
subgraph $D_*$ having all its induced original edges present and iid
with probability $p'$. Suppose each vertex $y$ in the collar satisfies
\[
 \PP_{p'}(B(y,m)\leftrightarrow T_*\text{ by paths in }D_*)
       \ge1-\eta_0.
\]
Suppose the wired seed vertices are outside the collar, and all other
edges of the product graph are original edges of $G$. Then
\begin{equation}\label{eq:fresh-collar-transfer}
 \PP(*\leftrightarrow B_*,\ *\nleftrightarrow T_*)\le\kappa.
\end{equation}
In particular this applies to each successive pair in
\eqref{eq:relay-regions}, whenever the whole corridor is fresh and
the wires are outside the input collar.
\end{proposition}

\begin{proof}
The empty input is immediate. We first choose a shell at which there
are unlikely to be only a few exterior contacts. We will use the
contacts to reach a relay set determined by independent interior bits,
then apply joint gluing after those bits have been fixed.

\paragraph{A shell with enough contacts.}
Use distance in the original graph to
$B_*$. For $m+1\le j\le R$, let $X_j$ be the level-$j$ vertices
reachable from $*$ using only vertices at levels at least $j$ and
the wired source. It is measurable in the exterior bits with both
endpoints at these levels. Hitting $B_*$ requires $X_j\ne\varnothing$.
If $1\le|X_j|<k$, closing all edges from $X_j$ to level $j-1$
has conditional probability at least $(1-\pc)^{Dk}$ and forces the
minimum reached level to be exactly $j$. These minimum-level events
are disjoint for different $j$. Thus some deterministic $j$ satisfies
\begin{equation}\label{eq:few-contacts}
 \PP(1\le|X_j|<k)
 \le\frac{(1-\pc)^{-Dk}}{R-m}<\kappa/4.
\end{equation}
The edges closed here are fresh collar edges, independent of the
exterior. Distance changes by at most one along every original edge;
the source wires lie outside all these levels.

Fix this deterministic $j$. Let $\mathcal F_{\rm ext}$ be the
sigma-field generated by edges with both endpoints at levels at least
$j$. It determines $X_j$ and will determine the finite list of contacts
used for trials.

\paragraph{An interior relay set.}
For every level-$j$ vertex $v$, fix an inward geodesic of length
$m+1$ ending at $y_v$, and a spanning tree of $B(y_v,m)$. The geodesic
and tree lie in the collar, and every one of their edges has an
endpoint of level less than $j$. Let $W_v$ be the event that all
these trial edges are open. Its probability satisfies
\[
 s_v:=\PP(W_v)\ge s_0:=(\pc/2)^{m+1+V_+(m)}.
\]
Let $\xi$ be the configuration on the union of all these fixed trial
edge sets, including trials at vertices not ultimately selected.
These are fresh product bits, independent of $\mathcal F_{\rm ext}$.
Define
\[
 g_v(\xi)=\PP(B(y_v,m)\leftrightarrow T_*\mid\xi).
\]
This conditional probability integrates over \emph{all} remaining
bits, including those generating $\mathcal F_{\rm ext}$. The actual
exterior chooses contacts but is not fixed in this target probability.
The function $g_v$ is increasing in $\xi$ and has
mean at least $1-\eta_0$. Harris's inequality and Markov's inequality
therefore give
\begin{equation}\label{eq:good-trial}
 \PP(W_v,\ g_v<1-\lambda)\le s_v\eta_0/\lambda.
\end{equation}

Define the relay set before choosing the contact list:
\[
 A(\xi)=\{v:\ d(v,B_*)=j,\ W_v\text{ occurs},\ g_v(\xi)\ge1-\lambda\}.
\]
This set contains every good trial vertex and depends only on $\xi$.
Given $\xi$, each $v\in A(\xi)$ is connected by its trial edges to
the whole seed ball, so its conditional probability of reaching $T_*$
is at least $1-\lambda$. The remaining law is a countable product
with the trial states fixed. Corollary~\ref{cor:gluing-failure} gives
\begin{equation*}
 \PP(*\leftrightarrow A(\xi),\ *\nleftrightarrow T_*\mid\xi)
       \le\lambda.
\end{equation*}
It remains to bound the probability that the source hits $B_*$
without reaching this relay set.

\paragraph{Reaching an interior relay from the exterior.}
Fix an enumeration of the countable vertex set. If $|X_j|\ge k$,
select its first available vertices in this order, using only the
exterior, at pairwise distances greater than $4m+2$, until
the sum of their $s_v$ is in $[K,K+1]$. Greedy selection deletes at
most $V_+(4m+2)$ candidates each time, so
\eqref{eq:corridor-trial-scales} guarantees this is possible; the
selected list is finite even when $X_j$ is infinite. Each trial lies
within distance $2m+1$ of its vertex, so the selected trials have
disjoint edge sets. Conditional on the exterior, their probability
of no success is at most $e^{-K}=\kappa/4$. By
\eqref{eq:good-trial} their probability of a successful trial with
$g_v<1-\lambda$ is at most
$(K+1)\eta_0/\lambda\le\kappa/4$. These estimates remain valid
after selection because the entire $\xi$ law is independent of
the exterior selecting the list.

On a hit of $B_*$, the contact set is nonempty. Outside the small-contact
error in \eqref{eq:few-contacts} and the two trial errors, some selected
exterior contact belongs to $A(\xi)$ and is already connected to the
source. Thus the probability of hitting $B_*$ without hitting $A(\xi)$
is at most $3\kappa/4$. Integrating the conditional gluing bound and
using $\lambda=\kappa/4$ proves \eqref{eq:fresh-collar-transfer}.
\end{proof}

\subsection{Adaptive exploration with stored predictions}

Stored high conditional probabilities are used in the proof of
Kozma and Nitzan's Theorem~6, particularly equations~(30)--(32)
\cite{KozmaNitzan2024}. The adaptive block strategy has an earlier
antecedent in Grimmett and Marstrand \cite[Sections~3--4]{GrimmettMarstrand1990}.
Here we use an oriented quadrant, choose one incoming corridor,
reveal whole edge domains, and prove the conditional-failure boundary
estimate explicitly.

\begin{proposition}[Subcritical exploration]\label{prop:exploration}
Suppose the coordinates and seed radius satisfy
Proposition~\ref{prop:moves}, the corridor layout satisfies
Lemma~\ref{lem:corridors}, and every relay in
\eqref{eq:relay-regions} satisfies Proposition~\ref{prop:transfer}.
If
\[
 2(\eta_0+L\kappa)<\delta,\qquad
 \delta+2(L+1)\kappa/\delta\le\varepsilon,
 \qquad\sum_{l\ge1}4^l\varepsilon^{l/4}<1,
\]
then Bernoulli bond percolation at $p'$ has an infinite cluster with
positive probability. In particular, the choices
\eqref{eq:corridor-trial-scales}--\eqref{eq:corridor-coordinate-step}
and the preceding error parameters have this consequence at $p'<\pc$.
\end{proposition}

\begin{proof}
Choose $o$ with $(t(o),z(o))=(0,0)$, and wire a synthetic source to
the finite ball $B(o,m)$. A revealed history
records a union of the edge domains $\mathcal Q_v,\mathcal E_{vx}$
and their actual $p'$ states. In a prospective graph obtained by adding
another edge domain, previously revealed states remain fixed and the
other edges have their conditional product law. We include only the
edges of those domains, even if further original edges join vertices
already present. The domains may be infinite; the conditional-law
justification below covers this point.

\paragraph{Initialization.}
At the root site, the two moves reach the corresponding $B_1$ directly
with probability at least $1-\eta_0$. From $B_1$ onward, source wires
are outside each input collar, by their time separation and
\eqref{eq:corridor-coordinate-step}. Summing
Proposition~\ref{prop:transfer} in a fixed child corridor shows that
the source misses its child mark with probability at most
$\eta_0+L\kappa$. Reveal $\mathcal Q_0$. The two conditional
predictions for reaching the child marks are both at least
$1-\delta$ on an event of positive probability: the probability
that either fails this threshold is at most
$2(\eta_0+L\kappa)/\delta<1$. Condition on this
$\mathcal Q_0$-measurable event, declare the root good, and store
these two predictions. The source already hits $M_0$ through $o$.
For the rest of the proof, probabilities refer to the law conditioned
on this initialization event.

\paragraph{A test and its two exposure times.}
Process sites by increasing layer $a+b$, in a fixed order within
each layer. A nonroot site is tested if it has a good parent; choose
one such parent $w$ deterministically. Sites with no good parent are
not tested. Before testing $v$, let $I$ denote the revealed history.
There are two geometric facts at this time. The incoming corridor
$\mathcal E_{wv}$ meets $I$ only in its parent's already known
$\mathcal Q_w$. Every outgoing corridor $\mathcal E_{vx}$ misses $I$
entirely. Indeed, no revealed corridor has endpoint $v$ or either child
of $v$: earlier tests reveal only their own incoming corridors, and
all sites of a preceding layer are processed first. Apply
Lemma~\ref{lem:corridors} to these domains and to $\mathcal Q_0$.

The first fact preserves the prediction stored when $w$ became good.
Intervening reveals meet $\mathcal E_{wv}$ only in $\mathcal Q_w$,
which was already known then, so they leave the other incoming bits
independent with parameter $p'$. The newly included history edges can
only add paths. Thus, in the prospective graph $I\cup\mathcal E_{wv}$,
\begin{equation}\label{eq:stored-prediction}
 \PP(A_v\mid I)\ge1-\delta,
 \qquad A_v=\{*\leftrightarrow M_v\text{ in }I\cup\mathcal E_{wv}\}.
\end{equation}

The second fact allows us to estimate both outgoing predictions
\emph{before} revealing the incoming corridor. For a child $x$ of $v$,
use the prospective graph
\[
 I\cup\mathcal E_{wv}\cup\mathcal E_{vx}.
\]
Conditionally on $I$, every edge of $D_{vx}$ still has its fresh iid
$p'$ law. In particular this includes the shared endpoint region
$\mathcal Q_v=\mathcal E_{wv}\cap\mathcal E_{vx}$, which has not yet
been revealed. The source wires lie outside all the input collars,
by time separation. Apply Proposition~\ref{prop:transfer} to the
$L+1$ consecutive relays, all in this same prospective graph. It gives
\begin{equation}\label{eq:prospective-transfer}
 \PP(A_v,\ *\nleftrightarrow M_x\mid I)\le(L+1)\kappa.
\end{equation}
A hit of the first region and a miss of the last entails a hit
followed by a miss at some consecutive pair, so no independence
between the successive estimates is needed.

Now reveal all of $\mathcal E_{wv}$ and call the enlarged history
$F$. The event $A_v$ and every state in $\mathcal Q_v$ are now known.
The outgoing edges outside $\mathcal Q_v$ remain unexposed. For each
child compute
\[
 q_x(F)=\PP(*\leftrightarrow M_x\text{ in }F\cup\mathcal E_{vx}\mid F).
\]
This product integral fixes the incoming states and integrates over
the unexposed outgoing edges; computing it reveals no further edge.
Declare $v$ good if $A_v$ occurs and both predictions are at least
$1-\delta$; otherwise declare the tested site failed. In the good
case store both predictions. By conditioning
\eqref{eq:prospective-transfer} on $F$,
\[
 \EE[\mathbf1_{A_v}(1-q_x(F))\mid I]\le(L+1)\kappa.
\]
Together with \eqref{eq:stored-prediction} and Markov's inequality,
this gives the uniform conditional failure bound
\begin{equation}\label{eq:conditional-test-failure}
 \PP(v\text{ fails}\mid I)
 \le\delta+\frac{2(L+1)\kappa}{\delta}
 \le\varepsilon.
\end{equation}

\paragraph{Conditional laws for countable reveal domains.}
We justify the product-law statements also for infinite reveal
domains. At any finite slot there are only finitely many earlier
decisions and parent choices. For each possible decision sequence,
the union $R$ of revealed edge sets is deterministic, and the event
of following that sequence is measurable in the coordinates of $R$.
The sequence, its edge states, and every prediction are functions
of those coordinates; a prediction is a measurable product integral
on the countable complement. On this part of the history space,
independence of disjoint coordinate sets gives the original product
law on $R^c$, conditionally on the full revealed states. This is a
statement about conditional kernels, and does not condition on a
positive-probability individual infinite configuration. It follows
either directly by Fubini on the two countable coordinate sets or by
first testing cylinder functions on $R^c$ and then using a monotone
class argument. The initial conditioning is measurable in
$\mathcal Q_0$ and hence preserves this property. This proves the
induction on slots. In particular, future corridor bits are not
revealed by decisions based on their conditional probabilities.

\paragraph{A boundary of failed sites.}
For any fixed finite set $S$ of nonroot sites,
\begin{equation}\label{eq:specified-failures}
 \PP(\text{every site of }S\text{ is tested and fails})
       \le\varepsilon^{|S|}.
\end{equation}
Indeed, enumerate the finite slots up to the last site of $S$.
At its slot, the event that a given site is tested and that the
previous specified sites failed is measurable before the test.
Apply \eqref{eq:conditional-test-failure} and iterate. A specified
site that is not tested makes the event impossible. A set containing
the good root also has probability zero.

If the good set is finite, draw the counterclockwise boundaries of
its unit site squares and cancel oppositely oriented shared edges.
The remaining directed graph is finite and balanced, so its edges
decompose into closed directed trails with no repeated edge. One
trail contains the fixed west edge of the root square, since the
sites lie in the nonnegative quadrant. In a trail of length $l$,
half the steps point up or left, by horizontal and vertical balance.
An up step crosses an east adjacency from a good site to a nongood
child; a left step crosses a north adjacency of the same kind.
Every such child was tested and failed, because all its parents were
processed earlier. Each child has at most two parents, so this trail
specifies at least $l/4$ distinct failed sites.

There are at most $4^l$ trails of length $l$ starting with the fixed
oriented root edge. For each candidate, its up and left edges
determine the outside child sites; restricting to edge-distinct
trails preserves the preceding multiplicity bound. Equation
\eqref{eq:specified-failures} and a union bound show that the
probability of a finite good set is at most the sum in
\eqref{eq:directed-contour-series}, which is less than one.

Thus with positive conditional probability there are infinitely
many good sites. Their layers, hence their time coordinates, are
unbounded. Each good mark was reached by an actual open path from
the wired finite seed at its test, and later reveals do not change
those states. The seed therefore has an infinite open component.
Removing the synthetic source leaves an infinite cluster from one
of its finitely many wired vertices. The initialization event had
positive probability, so this conclusion also has positive
probability in the unconditioned $p'$ law.

Finally, the chosen parameters satisfy the two numerical hypotheses.
They give $\eta_0\le\kappa$ and
\[
 2(\eta_0+L\kappa)\le2(L+1)\kappa
   =\varepsilon\delta/4<\delta,
 \qquad
 \delta+2(L+1)\kappa/\delta=\varepsilon/2\le\varepsilon.
\]
The remaining hypothesis is exactly
\eqref{eq:directed-contour-series}.
\end{proof}

Since $p'<\pc$, Proposition~\ref{prop:exploration} contradicts the
assumed critical jump and completes the induction in
Theorem~\ref{thm:structured}.
